\documentclass[UTF8,11pt]{ctexart}
\usepackage[a4paper,margin=24mm]{geometry}
\usepackage{amsmath,amssymb,amsthm,mathtools,mathrsfs}
\usepackage[all]{xy}
\usepackage{xurl,hyperref,enumitem}
\usepackage{tikz}
\usetikzlibrary{arrows.meta,cd,decorations.markings,calc,patterns}
\hypersetup{hidelinks,breaklinks=true}
\setlength{\emergencystretch}{3em}
\allowdisplaybreaks
\newtheorem*{theorem}{Theorem}
\newtheorem*{thm}{Theorem}
\newtheorem*{lemma}{Lemma}
\newtheorem*{proposition}{Proposition}
\newtheorem*{prop}{Proposition}
\newtheorem*{corollary}{Corollary}
\newtheorem*{cor}{Corollary}
\newtheorem*{claim}{Claim}
\theoremstyle{definition}
\newtheorem*{definition}{Definition}
\newtheorem*{defn}{Definition}
\newtheorem*{remark}{Remark}
\newtheorem*{example}{Example}
\newcommand{\R}{\mathbb R}
\newcommand{\C}{\mathbb C}
\newcommand{\Z}{\mathbb Z}
\newcommand{\Q}{\mathbb Q}
\newcommand{\N}{\mathbb N}
\newcommand{\F}{\mathbb F}
\newcommand{\D}{\mathbb D}
\newcommand{\bB}{\mathbb B}
\newcommand{\bC}{\mathbb C}
\newcommand{\bR}{\mathbb R}
\newcommand{\bZ}{\mathbb Z}
\newcommand{\bN}{\mathbb N}
\newcommand{\bQ}{\mathbb Q}
\newcommand{\bD}{\mathbb D}
\newcommand{\bF}{\mathbb F}
\newcommand{\CP}{\mathbb{CP}}
\newcommand{\RP}{\mathbb{RP}}
\newcommand{\sO}{\mathscr O}
\newcommand{\sI}{\mathscr I}
\newcommand{\sF}{\mathscr F}
\newcommand{\sG}{\mathscr G}
\newcommand{\sH}{\mathscr H}
\newcommand{\sR}{\mathscr R}
\newcommand{\sM}{\mathscr M}
\newcommand{\sV}{\mathscr V}
\newcommand{\sU}{\mathscr U}
\DeclareMathOperator{\Hom}{Hom}
\DeclareMathOperator{\Ext}{Ext}
\DeclareMathOperator{\Tor}{Tor}
\DeclareMathOperator{\Spec}{Spec}
\DeclareMathOperator{\Proj}{Proj}
\DeclareMathOperator{\supp}{supp}
\DeclareMathOperator{\im}{im}
\DeclareMathOperator{\id}{id}
\DeclareMathOperator{\Tr}{Tr}
\DeclareMathOperator{\tr}{tr}
\DeclareMathOperator{\rank}{rank}
\DeclareMathOperator{\ord}{ord}
\DeclareMathOperator{\dist}{dist}
\DeclareMathOperator{\End}{End}
\DeclareMathOperator{\Aut}{Aut}
\DeclareMathOperator{\Gal}{Gal}
\DeclareMathOperator{\vol}{vol}
\DeclareMathOperator{\Cl}{Cl}
\DeclareMathOperator{\coker}{coker}
\DeclareMathOperator{\codim}{codim}
\DeclareMathOperator{\divergence}{div}
\DeclareMathOperator{\Ric}{Ric}
\DeclareMathOperator{\grad}{grad}
\DeclareMathOperator{\Hess}{Hess}
\newcommand{\abs}[1]{\left\lvert#1\right\rvert}
\newcommand{\sabs}[1]{\lvert#1\rvert}
\newcommand{\babs}[1]{\bigl\lvert#1\bigr\rvert}
\newcommand{\Babs}[1]{\Bigl\lvert#1\Bigr\rvert}
\newcommand{\bbabs}[1]{\biggl\lvert#1\biggr\rvert}
\newcommand{\BBabs}[1]{\Biggl\lvert#1\Biggr\rvert}
\newcommand{\norm}[1]{\left\lVert#1\right\rVert}
\newcommand{\snorm}[1]{\lVert#1\rVert}
\newcommand{\bnorm}[1]{\bigl\lVert#1\bigr\rVert}
\newcommand{\Bnorm}[1]{\Bigl\lVert#1\Bigr\rVert}
\newcommand{\bbnorm}[1]{\biggl\lVert#1\biggr\rVert}
\newcommand{\BBnorm}[1]{\Biggl\lVert#1\Biggr\rVert}
\newcommand{\widebar}[1]{\overline{#1}}
\newcommand{\nicefrac}[2]{\frac{#1}{#2}}
\newcommand{\linnprod}[2]{\langle#1,#2\rangle}
\newcommand{\myindex}[1]{#1}
\newcommand{\glsadd}[1]{}
\newcommand{\avoidbreak}{}
\newcommand{\sourceref}[1]{\textnormal{[原文：\nolinkurl{#1}]}}
\newcommand{\thmref}[1]{\sourceref{#1}}
\newcommand{\lemmaref}[1]{\sourceref{#1}}
\newcommand{\propref}[1]{\sourceref{#1}}
\newcommand{\corref}[1]{\sourceref{#1}}
\newcommand{\defnref}[1]{\sourceref{#1}}
\newcommand{\exampleref}[1]{\sourceref{#1}}
\newcommand{\exerciseref}[1]{\sourceref{#1}}
\newcommand{\figureref}[1]{\sourceref{#1}}
\newcommand{\sectionref}[1]{\sourceref{#1}}
\newcommand{\appendixref}[1]{\sourceref{#1}}
\newcommand{\stacksref}[2]{\href{https://stacks.math.columbia.edu/tag/#1}{#2 [#1]}}

\begin{document}
\section*{036\quad Bergman再生核的构造与正交展开}
\subsection*{来源与计数目标}
Ji\v{r}\'i Lebl, \emph{Tasty Bits of Several Complex Variables}。从作者公开的原生 TeX 正文摘录；获取日期：2026-09-28。使用实际访问的在线正文，不把工作分支冒称冻结的印刷版。
\par\url{https://raw.githubusercontent.com/jirilebl/scv/master/scv.tex}
\par\url{https://www.jirka.org/scv/}
\par 原文位置：Chapter 5，\emph{The Bergman kernel}；\nolinkurl{lemma:bergmanKbound}、Riesz 构造段、共轭对称命题和正交展开命题。
\par\textbf{本文件只计一个问题：}一般复域上 Bergman 核的构造及正交展开，不把 Hilbert 空间完备性或对称性另计题目。
\subsection*{作者命题与人类证明}

Let $U\subset\C^n$ be a domain. Define the Bergman space of $U$:
\[A^2(U):=\sO(U)\cap L^2(U).\]
That is, $A^2(U)$ denotes the space of holomorphic functions $f\in\sO(U)$ such that
\[\|f\|_{A^2(U)}^2:=\|f\|_{L^2(U)}^2=\int_U|f(z)|^2dV<\infty.\]
The space $A^2(U)$ is an inner product space with the $L^2(U)$ inner product
\[\langle f,g\rangle:=\int_U f(z)\overline{g(z)}\,dV.\]
We will prove that $A^2(U)$ is complete; in other words, it is a Hilbert space. We first prove that the $A^2(U)$ norm bounds the uniform norm on compact sets.
\begin{lemma}\label{lemma:bergmanKbound}
Let $U\subset\C^n$ be a domain and $K\subset\subset U$ compact. Then there exists a constant $C_K$, such that
\[\|f\|_K=\sup_{z\in K}|f(z)|\leq C_K\|f\|_{A^2(U)}\qquad\text{for all }f\in A^2(U).\]
Consequently, $A^2(U)$ is complete.
\end{lemma}
\begin{proof}
As $K$ is compact there exists an $r>0$ such that $\overline{\Delta_r(z)}\subset U$ for all $z\in K$. Take any $z\in K$, and apply \exerciseref{exercise:averageDelta} and Cauchy--Schwarz:
\begin{align*}
|f(z)|&=\left|\frac1{V(\Delta_r(z))}\int_{\Delta_r(z)}f(\xi)\,dV(\xi)\right|\\
&\leq\frac1{\pi^nr^{2n}}\sqrt{\int_{\Delta_r(z)}1^2\,dV(\xi)}\sqrt{\int_{\Delta_r(z)}|f(\xi)|^2\,dV(\xi)}\\
&=\frac1{\pi^{n/2}r^n}\|f\|_{A^2(\Delta_r(z))}\leq\frac1{\pi^{n/2}r^n}\|f\|_{A^2(U)}.
\end{align*}
Taking the supremum over $z\in K$ proves the estimate. Therefore, if $\{f_\ell\}$ is a sequence of functions in $A^2(U)$ converging in $L^2(U)$ to some $f\in L^2(U)$, then it converges uniformly on compact sets, and so $f\in\sO(U)$. Consequently, $A^2(U)$ is a closed subspace of $L^2(U)$, and hence complete.
\end{proof}
The lemma says that point evaluation is a bounded linear functional. That is, for a fixed $z\in U$, taking $K=\{z\}$, the linear operator $f\mapsto f(z)$ is a bounded linear functional. By the Riesz representation theorem, there exists a $k_z\in A^2(U)$, such that
\[f(z)=\langle f,k_z\rangle.\]
Define the Bergman kernel for $U$ as
\[K_U(z,\bar\zeta):=\overline{k_z(\zeta)}.\]
The function $K_U$ is defined as $(z,\bar\zeta)$ vary over $U\times U^*$, where we write
\[U^*=\{\zeta\in\C^n:\bar\zeta\in U\}.\]
Then for all $f\in A^2(U)$, we have
\begin{equation}\label{eq:repropropBergman}
f(z)=\int_U f(\zeta)K_U(z,\bar\zeta)\,dV(\zeta).
\end{equation}
This last equation is sometimes called the reproducing property of the kernel. Note that the Bergman kernel depends on $U$, which is why we write it as $K_U(z,\bar\zeta)$.
\begin{prop}
The Bergman kernel $K_U(z,\bar\zeta)$ is holomorphic in $z$, antiholomorphic in $\zeta$, and
\[\overline{K_U(z,\bar\zeta)}=K_U(\zeta,\bar z).\]
\end{prop}
\begin{proof}
As each $k_z$ is in $A^2(U)$, it is holomorphic in $\zeta$. Hence, $K_U$ is antiholomorphic in $\zeta$. If we prove $\overline{K_U(z,\bar\zeta)}=K_U(\zeta,\bar z)$, then we prove $K_U$ is holomorphic in $z$. As $\overline{K_U(z,\bar\zeta)}=k_z(\zeta)$ is in $A^2(U)$, we have
\begin{align*}
\overline{K_U(z,\bar\zeta)}&=\int_U\overline{K_U(z,\bar w)}K_U(\zeta,\bar w)\,dV(w)\\
&=\overline{\left(\int_U\overline{K_U(\zeta,\bar w)}K_U(z,\bar w)\,dV(w)\right)}\\
&=\overline{\overline{K_U(\zeta,\bar z)}}=K_U(\zeta,\bar z).
\end{align*}
\end{proof}
Therefore, thinking of $\bar\zeta$ as the variable, $K_U$ is a holomorphic function of $2n$ variables.

\begin{prop}
Suppose $U\subset\C^n$ is a domain, and $\{\varphi_\ell(z)\}_{\ell\in I}$ is a complete orthonormal system for $A^2(U)$. Then
\[K_U(z,\bar\zeta)=\sum_{\ell\in I}\varphi_\ell(z)\overline{\varphi_\ell(\zeta)},\]
with uniform convergence on compact subsets of $U\times U^*$.
\end{prop}
\begin{proof}
For a fixed $\zeta\in U$, the function $z\mapsto K_U(z,\bar\zeta)$ is in $A^2(U)$. Expand this function in terms of the basis and use the reproducing property of $K_U$:
\begin{align*}
K_U(z,\bar\zeta)&=\sum_{\ell\in I}\left(\int_U K_U(w,\bar\zeta)\overline{\varphi_\ell(w)}\,dV(w)\right)\varphi_\ell(z)\\
&=\sum_{\ell\in I}\overline{\varphi_\ell(\zeta)}\varphi_\ell(z).
\end{align*}
The convergence is in $L^2$ as a function of $z$, for a fixed $\zeta$. Let $K\subset\subset U$ be a compact set. Via Lemma \ref{lemma:bergmanKbound}, $L^2$ convergence in $A^2(U)$ is uniform convergence on compact sets. Therefore, for a fixed $\zeta$ the convergence is uniform in $z\in K$. In particular, we get pointwise convergence. So,
\[\sum_{\ell\in I}|\varphi_\ell(z)|^2=\sum_{\ell\in I}\varphi_\ell(z)\overline{\varphi_\ell(z)}=K_U(z,\bar z)\leq C_K<\infty,\]
where $C_K$ is the supremum of $K_U(z,\bar\zeta)$ on $K\times K^*$. Hence for $(z,\bar\zeta)\in K\times K^*$,
\[
\sum_{\ell\in I}|\varphi_\ell(z)\overline{\varphi_\ell(\zeta)}|\leq\sqrt{\sum_{\ell\in I}|\varphi_\ell(z)|^2}\sqrt{\sum_{\ell\in I}|\varphi_\ell(\zeta)|^2}\leq C_K<\infty.
\]
So the convergence is uniform on $K\times K^*$.
\end{proof}

\subsection*{整理说明（非作者正文）}
完整收入从局部均值估计到点值泛函、Riesz 代表、共轭对称和正交展开的作者论证。略去穿插的圆盘显式计算、无界域例子和不参与该证明的练习。标准先修为多圆盘均值性质、Hilbert 空间 Riesz 表示及正交展开、全纯函数紧一致极限和 Cauchy--Schwarz；均值性质在原文以习题引用，本条明示其先修地位。保留原文末段关于紧集上一致收敛的论证，不作新增数学审稿或补造证明。

难度依据是将解析点值控制转成泛函表示，再恢复同时依赖两个复变量的再生核及基展开，而不是把 Riesz 定理直接套在未经确认的点值泛函上。正文原先无编号的命题仍用无编号环境；附加局部引用采用明确的标签文本。
\subsection*{可修改的简短标注（非作者正文）}
$c$：平方可积全纯函数空间的再生核。

$a$：多圆盘均值性质；局部一致估计；Hilbert 空间；有界点值泛函；Riesz 表示；共轭对称；正交展开。

\texttt{cross\_theory\_status = yes}。将全纯函数的局部解析估计转化为 Hilbert 空间有界泛函，再把代表向量和正交基展开还原成复解析再生核。位置：局部估计引理、$k_z$ 的构造和最后的正交展开。

全部标注均为非穷尽、可修改初稿；未标注不表示未使用。
\subsection*{许可与修改记录}
原数学正文作者：Ji\v{r}\'i Lebl。作者网站及源书声明允许 CC BY-SA 4.0 或 CC BY-NC-SA 4.0 双许可；本条选择 CC BY-SA 4.0，整理部分亦采用同一许可。\url{https://creativecommons.org/licenses/by-sa/4.0/}。2026-09-28：助手截取题证、调整独立排版，加入来源、范围与初稿标注；不暗示作者认可本次整理。保留的是所用 TeX 正文摘录，未将其称为整书原件。
\end{document}
